\documentclass[letter, 10pt]{article}
\usepackage[utf8]{inputenc}
\usepackage[T1]{fontenc}
\usepackage[spanish]{babel}
\usepackage{amsfonts}
\usepackage{amsmath}
\usepackage[dvips]{graphicx}
\usepackage{url}
\usepackage[top=3cm,bottom=3cm,left=3.5cm,right=3.5cm,footskip=1.5cm,headheight=1.5cm,headsep=.5cm,textheight=3cm]{geometry}



\begin{document}
\title{Inteligencia Artificial \\ \begin{Large}Estado del Arte: Problema MS-CFLP-CI\end{Large}}
\author{Benjamín Daza Jiménez}
\date{\today}
\maketitle


%--------------------No borrar esta secci\'on--------------------------------%
\section*{Evaluación}

\begin{tabular}{ll}
Resumen (5\%): & \underline{\hspace{2cm}} \\
Introducci\'on (5\%):  & \underline{\hspace{2cm}} \\
Definición del Problema (10\%):  & \underline{\hspace{2cm}} \\
Estado del Arte (40\%):  & \underline{\hspace{2cm}} \\
Modelo Matem\'atico (20\%): &  \underline{\hspace{2cm}}\\
Conclusiones (15\%): &  \underline{\hspace{2cm}}\\
Bibliograf\'ia (5\%): & \underline{\hspace{2cm}}\\
 &  \\
\textbf{Nota Final (100\%)}:   & \underline{\hspace{2cm}}
\end{tabular}
%---------------------------------------------------------------------------%
\vspace{2cm}


\begin{abstract}
La gestión de redes logísticas junto con restricciones de seguridad conlleva un compromiso fundamental entre la eficiencia de costos y la integridad operativa, especialmente bajo el modelo MS-CFLP-CI. La logística se vuelve un desafío que requiere diseñar sistemas que busquen minimizar gastos de apertura y transporte bajo segregaciones físicas estrictas entre los clientes en conflicto. El presente trabajo desarrolla una formulación matemática y evalúa algoritmos de vanguardia, con enfoque en la búsqueda de vecindarios grandes. Para ello, se establece que la flexibilidad del flujo fraccionado es el factor determinante para alcanzar la optimización en la industria, un entorno caracterizado por restricciones de capacidad y seguridad altamente restrictivas.
\end{abstract}

\section{Introducción}
Se ha vuelto una necesidad crítica para las industrias poder optimizar la localización de instalaciones bajo criterios de incompatibilidad entre clientes (MS-CFLP-CI), ya que manejan suministros reactivos, químicos o farmacéuticos. En estos sectores, una proximidad física poco adecuada no solo eleva costos operativos, sino que compromete la seguridad industrial y el cumplimiento normativo ambiental. Ante esto, se propone un análisis técnico que integra una formulación matemática con la revisión de las metaheurísticas más eficaces desarrolladas hasta la fecha. La contribución central es sintetizar criterios operativos para optimizar redes complejas bajo estándares de seguridad rigurosos.\\

La creciente necesidad de las industrias modernas de gestionar suministros de naturaleza crítica sustenta la motivación de este estudio, con productos tales como sustancias químicas reactivas, desechos industriales o productos farmacéuticos. La importancia de la investigación en estos sectores es máxima, ya que la gestión inadecuada de la proximidad física entre clientes que no son compatibles puede derivar en riesgos ambientales, accidentes operativos o incumplimientos normativos severos. Bajo este escenario, desarrollar modelos que garanticen la separación física mediante una distribución eficiente conforma un área de estudio vital para la seguridad industrial.\\

El problema a abordar consiste en determinar la configuración óptima de una red de distribución, decidiendo simultáneamente qué instalaciones deben ser abiertas o cerradas, y cómo se debe dividir la demanda de los clientes entre ellas. El principal desafío está en satisfacer los requerimientos de cada receptor, cumpliendo con las capacidades finitas de cada bodega y respetando las restricciones de incompatibilidad que prohíben que ciertos clientes puedan compartir un mismo centro de distribución.\\

El resto de este documento se organiza de la siguiente manera:

\begin{itemize}
    \item \textbf{Sección 2 - Definición del Problema:} Se describe la naturaleza del MS-CFLP-CI, detallando sus variables, restricciones operativas, objetivos económicos y la comparación con otras variantes y problemas relacionados.
    \item \textbf{Sección 3 - Estado del Arte:} Se presenta una revisión histórica de la evolución del problema, los algoritmos más eficaces desarrollados hasta la actualidad y las tendencias recientes en la disciplina.
    \item \textbf{Sección 4 - Modelo Matemático:} Se expone la formulación de programación lineal entera mixta (MIP), definiendo cada componente del modelo de optimización.
    \item \textbf{Sección 5 - Conclusiones:} Se sintetizan los hallazgos principales y se proponen futuras líneas de investigación.
\end{itemize}

\section{Definición del Problema}

\subsection{Descripción general}
El Capacitated Facility Location Problem with Customer Incompatibilities (CFLP-CI) se define como un problema de optimización combinatoria que extiende el modelo clásico de localización de instalaciones con capacidad limitada, el cual consiste en el diseño de una red de distribución que determine qué instalaciones abrir o cerrar, y cómo asignar la demanda de los clientes, cumpliendo con restricciones operativas y de seguridad \cite{ceschia2024}. En este estudio se abordará la variante Multi-Source (MS-CFLP-CI), permitiendo que múltiples bodegas abiertas puedan satisfacer la demanda de un mismo cliente de forma fragmentada; esto otorga mayor flexibilidad operativa frente a límites de capacidad ajustados \cite{Maia2023}.\\

La existencia de este modelo en la planificación logística es fundamental, en especial para industrias que gestionan productos cuya convivencia física es inviable, tales como sustancias químicas reactivas, desechos tóxicos o suministros alimenticios junto a pesticidas. En consecuencia, es indispensable considerar la separación física rigurosa, ya que no solo responde a eficiencia en costos, sino que son requisitos críticos para mitigar riesgos operativos, evitar contaminaciones cruzadas y cumplir con estándares de seguridad industrial junto con normativas legales \cite{Gjergji2025}.

\subsection{Dificultades y Complejidad del problema}
Abordar el MS-CFLP-CI significa considerar que colapsa rápidamente la eficiencia de los métodos exactos debido al grafo de conflictos . El problema base de localización se clasifica como NP-duro, lo que significa que no se conocen algoritmos capaces de encontrar la solución óptima en tiempo polinomial a medida que el número de instalaciones y clientes crece . En la práctica, la alta cantidad de restricciones de incompatibilidad convierte la búsqueda de una solución factible en el primer obstáculo, incluso antes de optimizar la estructura de los costos .\\

Más en específico, al expandir el espacio de búsqueda la complejidad de la variante multi-source se ve incrementada. A diferencia de modelos de fuente única, el sistema no solo enfrenta la decisión de apertura y asignación, sino que debe determinar de forma óptima la proporción de demanda que cada bodega debe suministrar a cada cliente \cite{Maia2023}. También se debe considerar que, para transicionar de un problema de asignación pura a uno de distribución de flujos fraccionados, se requiere de algoritmos robustos que puedan gestionar la importante cantidad de combinaciones posibles en redes de gran escala \cite{Gjergji2025}.

\subsection{Variables involucradas}
Para modelar el comportamiento de esta red logística y determinar la configuración óptima del sistema, el modelo toma decisiones basándose en dos tipos de variables fundamentales, descritas a continuación \cite{basu2014}:

\begin{itemize}
    \item \textbf{Variables de estado o apertura:} Operan como decisiones lógicas (activado/desactivado) que deciden si una bodega potencial debe ser abierta o permanecer cerrada. Es sumamente importante en el sistema porque habilita la capacidad operativa de una instalación y gatilla en el cobro de los costos fijos asociados a la infraestructura y mantenimiento.
    \item \textbf{Variables de flujo o asignación:} En la variante multi-source, representan la cantidad específica o fracción de la demanda que se despachará desde cada bodega activa hacia cada cliente. El comportamiento de estas variables es fundamental, ya que determina la distribución dinámica de la carga en la red, permitiendo que un cliente sea abastecido por múltiples fuentes, donde el modelo podrá optimizar el uso de las capacidades finitas de las bodegas mientras acumula de forma proporcional los costos variables de envío.
\end{itemize}

\subsection{Restricciones del problema}
El conjunto de soluciones factibles para este problema depende del cumplimiento estricto de las siguientes condiciones \cite{Maia2023}:

\begin{itemize}
    \item \textbf{Capacidad limitada:} Ninguna bodega puede procesar, almacenar o despachar un volumen de productos que exceda su límite de recursos predefinido. Funciona como un tope operativo que obliga al modelo a distribuir la carga de trabajo entre las distintas instalaciones activas de manera equilibrada.
    \item \textbf{Satisfacción de la demanda:} El requerimiento total de cada cliente debe ser cubierto por el sistema. Por la variante multi-source, se permite explícitamente que esta demanda sea satisfecha a partir de la suma de flujos provenientes de múltiples instalaciones, otorgando flexibilidad para cumplir con los pedidos.
    \item \textbf{Incompatibilidad de clientes:} Prohíbe que determinados pares de clientes, identificados previamente por conflictos normativos o de seguridad, sean abastecidos por la misma instalación física simultáneamente. Funciona como un filtro de seguridad que limita las combinaciones de asignación posibles dentro de una misma bodega \cite{ceschia2024, Gjergji2025}.
\end{itemize}

\subsection{Objetivos del problema}
El principal propósito de este modelo es optimizar la estructura de la red logística a través de minimizar los costos totales del sistema. El objetivo se alcanza al encontrar el equilibrio económico entre los dos componentes de costo principales, donde su priorización se puede analizar mediante reglas de decisión multicriterio con el fin de balancear las inversiones y gastos operativos \cite{Maia2023, Herrero2021}:

\begin{itemize}
\item \textbf{Costos de inversión inicial:} Gastos fijos requeridos para la apertura, habilitación y mantenimiento operativo de las bodegas seleccionadas dentro del conjunto de ubicaciones potenciales. Estos costos se incurren únicamente por activar una instalación.
\item \textbf{Costos de operación y transporte:} Gastos variables generados al satisfacer la demanda de cada cliente desde las bodegas activas. En la variante multi-source, estos se calculan de manera precisa en función de las fracciones de demanda enviadas y los costos unitarios de transporte asociados a cada par bodega-cliente.
\end{itemize}

\subsection{Variantes y problemas relacionados}
El Capacitated Facility Location Problem with Customer Incompatibilities (CFLP-CI) se deriva del problema clásico de localización de instalaciones (FLP) y su extensión capacitada (CFLP), siendo bases lógicas para la asignación de demandas que consideran los límites operativos en las infraestructuras \cite{Gjergji2025, Maia2023}. A partir de este modelo base, se distinguen dos variantes principales según la flexibilidad de abastecimiento permitida:

\begin{itemize}
    \item \textbf{Variante de abastecimiento limitado ($k$-source):} Algunos estudios recientes introducen una variante intermedia donde se permite el flujo múltiple, pero además, existe un límite superior $k$ sobre la cantidad de instalaciones que pueden servir a un mismo cliente (comúnmente $k=2$). Esta variante busca reducir la complejidad logística y los costos administrativos derivados de gestionar demasiados proveedores para un solo receptor \cite{ceschia2024}.
    \item \textbf{Problema de Ruteo de Vehículos con Cargas Incompatibles (VRP-IC):} Extiende el problema de ruteo vehicular al incorporar restricciones operativas donde ciertos productos o mercancías no pueden ser transportados juntos en el mismo vehículo. Esto responde a normativas de seguridad o a incompatibilades físicas. Esta variante busca minimizar costos de ruta respetando capacidad finita y la exclusión mutua de la carga \cite{pollaris2015}.
    \item \textbf{Problema de Asignación Generalizada con Conflictos (GAP-C):} Generaliza los modelos de asignación tradicionales mediante la introducción de un grafo de conflictos. Determinadas asignaciones son mutuamente excluyentes y no pueden ser procesadas simultáneamente por la misma instalación. Esta variante busca encontrar una asignación de costo mínimo que garantice la factibilidad frente a este conjunto de restricciones duras \cite{dambrosio2020}.
    \item \textbf{Clustering con Restricciones (Constrained K-Medoids):} En la minería de datos y localización de instalaciones, esta aproximación incorpora restricciones espaciales del tipo cannot-link. Esta condición obliga a que pares específicos de nodos o clientes en conflicto sean estrictamente asignados a distintos medoides (centros o bodegas), garantizando la separación de elementos incompatibles dentro de la topología de la red \cite{wagstaff2001}.
\end{itemize}

Adicionalmente, este problema tiene relación con el Hub Location Problem. Sin embargo, este introduce una jerarquía estructural distinta; mientras que el MS-CFLP-CI se enfoca en envíos directos de bodegas a clientes, el modelo de Hubs añade nodos de consolidación y transferencia intermedios, centralizando el tráfico y modificando radicalmente la dinámica de los costos de transporte y rutas de distribución \cite{Gjergji2025}.

\section{Estado del Arte}

\subsection{Origen y Evolución del Problema}
El estudio del CFLP es un área consolidada en la investigación de operaciones, pero la integración de incompatibilidades entre los clientes responde a una necesidad de modelar escenarios críticos, como la gestión de residuos industriales peligrosos o el almacenamiento de sustancias químicas reactivas \cite{Maia2023, owen1998}. Ante esto, considerar la separación física no es opcional, se transforma en un requisito de seguridad que las variantes tradicionales no lograban capturar. El modelo \textit{multi-source} ha ganado relevancia reciente, demostrando que la flexibilidad de flujos fraccionados permite considerar restricciones de seguridad sin sacrificar la viabilidad económica del sistema \cite{ceschia2024, Gjergji2025}.

\subsection{Métodos y Algoritmos de Resolución}
Se ha abordado el problema mediante una jerarquía de métodos, evolucionando en complejidad y eficacia:

\begin{itemize}
    \item \textbf{Métodos Completos:} Se ha utilizado Programación Lineal Entera Mixta (MIP) mediante solvers como Gurobi. Si bien garantizan optimización, su uso se limita a instancias pequeñas, ya que el grafo de conflictos aumenta exponencialmente el tiempo de cómputo en instancias reales \cite{Gjergji2025}.
    \item \textbf{Heurísticas y Metaheurísticas Iniciales:} Maia et al. \cite{Maia2023} introdujeron métodos como Iterated Local Search (ILS) y Algoritmos Evolutivos (PcEA). Se destacó MineReduce, una técnica que utiliza minería de datos para identificar patrones frecuentes de asignación y reducir el espacio de búsqueda.
    \item \textbf{Enfoques de Trayectoria:} Ceschia y Schaerf \cite{ceschia2024} propusieron un algoritmo de Recocido Simulado, el cual utiliza movimientos de intercambio y reasignación, logrando superar a las heurísticas previas gracias a un pre-procesamiento efectivo del grafo de incompatibilidades.
    \item \textbf{Algoritmo de Vanguardia (LNS):} Propuesta por Gjergji et al. \cite{Gjergji2025}. Este método utiliza operadores de destrucción (que eliminan parte de la solución actual) y reparación (que reconstruyen la solución de forma óptima), logrando los mejores resultados conocidos hasta la fecha.
\end{itemize}

\subsection{Representaciones y Movimientos}
La efectividad de los algoritmos ha dependido en gran medida de la representación de los datos y los tipos de movimientos empleados:

\begin{itemize}
    \item \textbf{Representaciones:} Las representaciones basadas en permutaciones han demostrado ser ineficientes, ya que los operadores de cruce suelen generar soluciones que no cumplen con las restricciones de incompatibilidad. Por otro lado, las representaciones basadas en listas de adyacencia para el grafo de conflictos y matrices de flujo para las asignaciones han permitido validaciones de factibilidad en tiempo casi constante \cite{Maia2023, Gjergji2025}.
    \item \textbf{Tipos de Movimientos:} Los algoritmos exitosos emplean movimientos de Reinsert, Swap y, en el caso de LNS, operadores de destrucción aleatoria o basada en costos (como eliminar los clientes con mayores costos de transporte) para escapar de óptimos locales.
\end{itemize}

\subsection{Tendencias Actuales}
La actual investigación está orientada hacia las Mateheurísticas, utilizando solvers exactos como Gurobi para resolver subproblemas de flujo de forma óptima durante las fases de reparación de las metaheurísticas, reduciendo el tiempo de cómputo en redes densas \cite{Gjergji2025}. También, el desarrollo de Adaptive LNS (ALNS) busca que el algoritmo pueda seleccionar dinámicamente operadores de destrucción, los cuales prioricen eliminar nodos que saturan el grafo de conflictos, logrando optimizar la exploración de regiones del espacio de búsqueda que los métodos estáticos normalmente ignoran \cite{ceschia2024}.

\section{Modelo Matemático}
El problema MS-CFLP-CI se formaliza mediante un modelo de Programación Lineal Entera Mixta (MIP). Esta formulación permite capturar la naturaleza continua de los flujos de demanda y la naturaleza discreta de las decisiones de apertura e incompatibilidad \cite{ceschia2024}.

\subsection{Formulación del Modelo}

\textbf{Conjuntos y Parámetros:}
\begin{itemize}
    \item $N$: Conjunto de clientes indexados por $i$.
    \item $M$: Conjunto de instalaciones potenciales indexadas por $j$.
    \item $I$: Conjunto de pares de clientes $\langle i, k \rangle$ que presentan incompatibilidad entre sí.
    \item $d_i$: Demanda total requerida por el cliente $i$.
    \item $s_j$: Capacidad de suministro máxima de la instalación $j$.
    \item $f_j$: Costo fijo incurrido al abrir y operar la instalación $j$.
    \item $c_{ij}$: Costo unitario de transporte para enviar productos desde la instalación $j$ al cliente $i$.
\end{itemize}

\textbf{Variables de Decisión:}
\begin{itemize}
    \item $y_j \in \{0, 1\}$: Variable binaria que toma el valor 1 si la instalación $j$ es abierta; 0 en caso contrario.
    \item $w_{ij} \in \{0, 1\}$: Variable binaria de asignación que indica si el cliente $i$ es abastecido (total o parcialmente) por la instalación $j$.
    \item $x_{ij} \in [0, 1]$: Variable continua que representa la fracción de la demanda del cliente $i$ satisfecha por la instalación $j$.
\end{itemize}

\textbf{Función Objetivo:}
\begin{equation}
    \min Z = \sum_{j \in M} f_j y_j + \sum_{i \in N} \sum_{j \in M} (c_{ij} \cdot d_i) x_{ij}
\end{equation}

\textit{Descripción:} El objetivo es minimizar la suma total de los costos. El primer término cuantifica los costos fijos de apertura de bodegas, mientras que el segundo término representa los costos variables de transporte, calculados sobre la cantidad real de producto despachado ($d_i x_{ij}$).\\

\textbf{Restricciones:}

1. \textbf{Satisfacción de Demanda:} 
\begin{equation}
    \sum_{j \in M} x_{ij} = 1, \quad \forall i \in N
\end{equation}

Define la cobertura total de los requerimientos para cada nodo receptor en la red.\\

2. \textbf{Límite de Capacidad y Activación:} 
\begin{equation}
    \sum_{i \in N} d_i x_{ij} \leq s_j y_j, \quad \forall j \in M
\end{equation}

Condiciona el volumen de flujo saliente a la capacidad nominal $s_j$ y restringe la factibilidad de despacho al estado de apertura de la instalación ($y_j$).\\

3. \textbf{Vínculo de Asignación Lógica:} 
\begin{equation}
    x_{ij} \leq w_{ij}, \quad \forall i \in N, \forall j \in M
\end{equation}

Impone la existencia de una asignación lógica ($w_{ij}=1$) como requisito estricto para habilitar flujos fraccionados positivos.\\

4. \textbf{Incompatibilidad de Clientes:} 
\begin{equation}
    w_{ij} + w_{kj} \leq 1, \quad \forall \langle i, k \rangle \in I, \forall j \in M
\end{equation}

Restricción de exclusión mutua que prohíbe la asignación simultánea de clientes en conflicto a un mismo centro logístico.\\

\section{Representación}

\section{Descripción del algoritmo}

\section{Experimentos}

\section{Resultados}

\section{Conclusiones}
El análisis del MS-CFLP-CI revela que se hace inviable la búsqueda de un óptimo global mediante métodos exactos a escala industrial, debido a la complejidad computacional que impone el grafo de conflictos \cite{ceschia2024}. Dada esta limitación, LNS se ha consolidado como la respuesta más efectiva para evadir las infactibilidades. El éxito de este enfoque se debe en gran medida a su capacidad para destruir y reconstruir dinámicamente la red, permitiendo explorar regiones del espacio de soluciones donde los algoritmos tradicionales colapsan \cite{Gjergji2025}.

\subsection{Análisis comparativo de las técnicas}
Luego de contrastar las metodologías presentes en el estado del arte, es evidente que los métodos exactos (MIP) quedan apartados con un rol referencial, sufriendo de una explosión combinatoria inmediata al escalar la densidad de las incompatibilidades. Por otro lado, metaheurísticas de trayectoria como Simulated Annealing logran solo factibilidad inicial, ya que demuestran ser ciegas al quedar atrapadas en óptimos locales profundos dictados por las restricciones de seguridad.\\

Actualmente, la Large Neighborhood Search (LNS) domina los resultados empíricos debido a su mecánica de reconstrucción parcial. Sin embargo, este enfoque presenta una debilidad estructural: depende de operadores de destrucción diseñados de forma manual o estática. Esto implica que, a pesar del éxito empírico de LNS, la literatura aún necesita de un método verdaderamente inteligente que logre aprender la topología del conflicto en tiempo real para destruir y reconstruir la red de forma autónoma.

\subsection{Limitaciones y tendencias}
A pesar de los avances, las técnicas actuales presentan limitaciones en escenarios de extrema saturación, donde las restricciones de capacidad y de incompatibilidad coinciden de tal manera que el espacio de soluciones factibles se vuelve casi nulo. En estos puntos de transición de fase, la mayoría de las metaheurísticas puras fallan al no contar con mecanismos de exploración lo suficientemente profundos. La tendencia actual sugiere que las mejores soluciones se logran mediante la combinación de exploración global (metaheurística) y refinamiento local exacto (solver matemático), permitiendo que el sistema escape de óptimos locales con mayor eficiencia.

\subsection{Lineamientos para investigaciones futuras}
El trabajo futuro en el área del MS-CFLP-CI ofrece oportunidades significativas para la innovación. Se proponen como lineamientos estratégicos para continuar la investigación:
\begin{itemize}
    \item \textbf{Implementación de operadores adaptativos:} El uso de \textit{Adaptive Large Neighborhood Search} (ALNS) para que el algoritmo aprenda, en tiempo real, qué operadores de destrucción son más efectivos según la morfología del grafo de conflictos de la instancia.
    \item \textbf{Incluir Aprendizaje Automático:} Se propone como idea original el uso de redes neuronales para predecir qué instalaciones tienen mayor probabilidad de pertenecer a la solución óptima, permitiendo un pre-procesamiento inteligente que reduzca drásticamente el espacio de búsqueda binario.
    \item \textbf{Modelamiento de incertidumbre:} Extender el estudio hacia variantes estocásticas donde la demanda o las capacidades sean variables, alineando el modelo matemático con la volatilidad real de las cadenas de suministro modernas.
\end{itemize}

\section{Uso LLMs}

Para el desarrollo del presente informe se utilizó una herramienta de Inteligencia Artificial (LLMs) con las siguientes condiciones:

\begin{enumerate}
    \item \textbf{Herramienta utilizada:} Gemini 3.1 (Pro)
    \item \textbf{Propósito de uso:} Corrección de faltas ortográficas, sugerencias de redacción, resumen de papers y, orientación en las limitaciones y tendencias.
    \item \textbf{Prompts relevantes:}
    \begin{itemize}
        \item Para el siguiente párrafo debes corregir faltas ortográficas y de puntuación: (párrafo).
        \item Para el siguiente párrafo debes sugerir redacción formal, adecuadas para un informe de investigación y considerando conectores usados con antelación: (párrafo).
        \item Para el siguiente paper, hazme un resumen acerca de (tema): (paper adjunto).
        \item Considerando lo descrito hasta el momento, que limitancias relevantes consideras.
    \end{itemize}
    \item \textbf{Nivel de edición:} el texto fue reescrito completamente, en algunos casos se tomó la sugerencia de redacción (solo frases). La información resumida fue complementada con el propio material adjunto.
    \item \textbf{Validación:} Cada resumen y sugerencia en las limitaciones y tendencias fueron validadas manualmente con el material del curso, los artículos de referencia y los papers propuestos.
\end{enumerate}

\section{Bibliografía}

\bibliographystyle{plain}
\bibliography{Referencias}

\end{document} 